\begin{abstract}
Textual-program optimization must balance proposing and evaluating candidates
under a fixed rollout budget.  Evaluating each candidate more extensively
leaves fewer rollouts for proposing new candidates.  We introduce \method{},
which turns accumulated evaluation results into relative credit to guide
subsequent proposals.  On each
shared instance, attaining the highest observed score earns more credit when a
smaller fraction of the other candidates attain it.  \method{} averages these
per-instance credits to select candidates for further revision and updates the
credits as new results arrive.  For revision, it assigns batches of instances
using the selected candidates' score gaps to the best score recorded among them
on each instance.  Each new candidate is evaluated on instances that were not
used to generate it or its ancestors.  Across six benchmarks, \method{}
achieves the highest average score among the compared methods for both
Qwen3-8B and GPT-5.4-nano.  On Qwen3-8B, it improves over GEPA by $11.1$
percentage points on average and $31.3$ points on AIME-2025.  Ablations on
Qwen3-8B show that both relative credit and matching improve the average score
across the six benchmarks.
\end{abstract}

\section{Introduction}
\label{sec:introduction}

\begin{wrapfigure}[21]{r}{0.47\textwidth}
\centering
\includegraphics[width=\linewidth]{figures/qwen_radar/qwen_radar.pdf}
\caption{\textbf{COMPASS leads on all six tasks.} \mbox{Qwen3-8B} scores (\%) use
labeled center-to-rim ranges (Table~\ref{tab:qwen_results}).}
\label{fig:qwen_radar}
\end{wrapfigure}

Feedback-driven optimizers revise instructions, prompts, and agent skills from
execution traces.  Reflection, textual gradients, and evolutionary prompt
search use language feedback to produce large behavioral updates with fixed
model weights \citep{reflexion,self_refine,textgrad,
cheng2024traceautodiffgenerativeoptimization,gepa}.  Given a seed program, a
benchmark evaluator, and rollout budget $B$, an optimizer generates and
evaluates candidates, then returns one program.  Rollouts provide traces for
proposing revisions and scores for comparing candidates.  Admission uses
a separate batch to decide whether a fixed revision enters the candidate pool.
For example, three instruction-following cases may expose a constraint-ordering
error and guide a new rule.  The rule is then evaluated on other cases.

Each proposal-and-evaluation attempt spends rollout budget on traces and
admission scores.
An observation can also contribute to later candidate comparisons and proposal
decisions.  Tracking which examples shaped each candidate determines how that
observation can be reused along the lineage.

The central question is how to use a growing, uneven record of execution
results to choose which candidates to revise and which batches to use next.
We call this \emph{provenance-constrained evidence routing}.  Proposal history
determines which observations can evaluate a candidate, and the recorded
results guide the next search decisions.  The aim is to maximize the final
program's expected benchmark utility under the available rollout budget.

This requires three linked decisions.  First, a candidate's proposal history
must follow its descendants when choosing admission examples.  Second, scores
collected on different subsets need to be compared on shared instances.  The
\emph{frontier} on an instance consists of its highest-scoring eligible
candidates and changes as results arrive.  Third, concurrent proposal attempts
need batches matched to each parent's recorded score gaps.

\method{} (\textbf{C}redit-Aware \textbf{O}ptimization with \textbf{M}ini-batch
\textbf{P}roposal and \textbf{A}dmission for \textbf{S}kill \textbf{S}earch)
links these decisions.  Recursive exclusion separates proposal and admission
examples throughout the lineage.  Relative credit updates candidate rankings
as instance frontiers change.  COMPASS selects the highest-credit candidates
for revision, including cutoff ties.  Maximum-weight
matching assigns batches to their proposal slots, each representing one
 revision attempt. Each fixed child receives a
separately sampled legal admission batch.  Figure~\ref{fig:overview} follows
this process from candidate generation to final selection.

\begin{figure}[t]
\centering
\includegraphics[width=\linewidth]{figures/figure1/figure1_compass_evidence_routing.pdf}
\caption{\textbf{Candidate comparisons guide the next search wave in \method{}.}
Child $p$, proposed from parent $s$ on $P$, is evaluated on a separate legal
batch $A$ and enters the cache after admission.  Legal same-instance comparisons
produce credit $G_s$, which selects proposal parents, admission references, and
the final program.  $H$ matches parent slots to proposal batches, while each
admission batch is sampled independently outside $X(s,P)$.  Red hatching marks
lineage-excluded observations.}
\label{fig:overview}
\end{figure}

Completed benchmark evaluations span model families and task types.  On
Qwen3-8B, the programs selected by COMPASS for each task reach $63.33\%$ on AIME-2025,
$43.37\%$ on IFBench, $73.96\%$ on LiveBench-Math, and $54.67\%$ on HoVer.
These scores exceed published GEPA by $31.33$, $4.76$, $22.01$, and $2.34$
percentage points, respectively.

This paper makes four contributions.
\begin{enumerate}[leftmargin=*,itemsep=2pt,topsep=3pt]
  \item \emph{Provenance-constrained evidence routing} formulates fixed-budget
  search and final selection with proposal history tracked throughout each
  lineage (\Cref{sec:setting}).

  \item \method{} converts eligible same-instance comparisons into relative
  credit and updates that credit when new results arrive
  (\Cref{sec:global_credit,sec:lineage_compression}).

  \item Relative credit selects proposal parents, and maximum-weight matching
  assigns their batches using recorded score gaps.  Paired admission evaluates
  each resulting revision (\Cref{sec:proposal_allocation}).

  \item Evaluations under fixed published budgets cover two models and six
  tasks.  COMPASS has the highest six-task aggregate in both model comparisons
  (\Cref{sec:experiments}).
\end{enumerate}

\section{Provenance-Constrained Evidence Routing}
\label{sec:setting}

\subsection{Budgeted textual-program optimization}

Given a seed textual program $p_0$, a benchmark evaluator, and a rollout
budget $B$, a complete run $\pi$ specifies the proposal and admission decisions
made within that budget.  Each attempt in the run uses a proposal batch $P$ of
$b_P$ instances to generate a candidate and a separate admission batch $A$ of
$b_A$ instances to evaluate it.  Proposal generation and admission evaluation
share the rollout budget.  COMPASS starts a new wave while the cumulative
proposal and evaluation cost remains below $B$, then completes that wave before
stopping.  Every generated child is evaluated on $b_A$ admission instances.

Let $\Pi(B,b_P,b_A)$ be the set of complete runs with proposal batch size
$b_P$, admission batch size $b_A$, and budget $B$.  Let
$p_{\mathrm{out}}(\pi)$ denote the final program returned by run $\pi$.
Let $R(p)$ denote the expected benchmark utility of program $p$, and let
$R^\star$ denote the best utility reachable under the search model.  The
budgeted optimization problem is
\[
\pi^\star\in
\arg\min_{\pi\in\Pi(B,b_P,b_A)}
\mathbb{E}\!\left[
R^\star-R\!\left(p_{\mathrm{out}}(\pi)\right)
\right].
\]

At optimization step $t$, let $\Pcal_t=\{p_0,\ldots,p_{m_t-1}\}$ contain the
accepted textual programs.  Every candidate after $p_0$ records a parent,
inducing a directed acyclic lineage.  The \emph{benchmark evaluator} runs each
program using the benchmark's official data, execution procedure, scoring
metric, and aggregation rule.
Executing program $p_s$ on optimization instance $x$ yields a bounded reward in
$[0,1]$, an output, and an execution trace.  The committed cache is a partial
table $q_t(s,x)$ containing the highest valid, finite score observed for each
evaluated pair, with unevaluated pairs represented by missing entries.  A new
observation replaces an existing score and its associated output only when its
score is strictly higher.  Thus each candidate--instance pair retains one
record.  Let
$P_a^{\mathrm{birth}}$ be the proposal batch that created candidate $a$, with
$P_0^{\mathrm{birth}}=\varnothing$, and let $\operatorname{Anc}(s)$ contain the
transitive ancestors of $s$.  The proposal history of candidate $s$ is
\[
\mathcal B_s=\bigcup_{a\in\operatorname{Anc}(s)\cup\{s\}}
P_a^{\mathrm{birth}}
\]
The credit-eligible evidence set is
$\Omega_t(s)=\{x\mid q_t(s,x)\text{ exists and }x\notin\mathcal B_s\}$.
Relative credit for candidate $s$ therefore draws exclusively from observations
outside the batches that shaped its lineage.

The benchmark evaluator supplies the optimization observations stored in the cache and,
after selection, an independent frozen test.

\subsection{Separating proposal and admission evidence}

For a child proposed from parent $s$ on batch $P$, its proposal history combines
$\mathcal B_s$ with the current batch.
\begin{equation}
X(s,P)=P\cup
\bigcup_{a\in\operatorname{Anc}(s)\cup\{s\}}P_a^{\mathrm{birth}}
\label{eq:recursive_exclusion}
\end{equation}
A legal admission batch satisfies $A\cap X(s,P)=\varnothing$.  For example,
suppose the three diagnostic cases help produce $p_1$ from $p_0$, and $p_2$ is
later revised from $p_1$ using another batch.  The original cases remain in
$p_2$'s inherited proposal history.  Admission for both $p_1$ and $p_2$ therefore
draws outside the history that shaped each candidate, while the proposal traces
remain available for diagnosis.

\subsection{Choosing the next search actions}

Let $\mathcal H_t$ contain the lineage, proposal batches, committed cache,
remaining budget, and persistent routing state.  A routing action
$a_t=(\A_t,\mathbf z_t,Y_t)$ chooses active candidates, proposal parents,
and a one-to-one assignment of their proposal slots to batches.
Relative credit selects candidates, and matching assigns batches using their
recorded score gaps.  Each fixed child receives a separately sampled admission
batch outside its proposal history.  The frozen benchmark test is reserved for
final evaluation.

A feasible COMPASS run enforces this exclusion, the wave-level stopping rule,
one cached record per candidate--instance pair, and a separate final test.  The
best such run minimizes expected simple regret for the returned program.

\section{Candidate Selection, Matching, and Admission in \method{}}
\label{sec:method}

Each wave selects proposal parents and assigns them batches using the current
cache.  The resulting traces guide revisions, which are evaluated on legal
admission batches with fixed references.  Accepted outcomes update the
comparison state and guide the next wave.  The following sections define the
credit, parent selection, matching, and admission steps.

\subsection{Relative credit from shared observations}
\label{sec:global_credit}

Per-instance credit captures distinctions among candidates on the same instance.  For
example, if three candidates all succeed or all fail on a binary-scored
instance, each receives zero credit.  A single success gives credits
$(1,-\tfrac12,-\tfrac12)$. Two successes give
$(\tfrac12,\tfrac12,-1)$, with successful candidates listed first.  The
following definition extends this comparison to every legally observed frontier.

Let $\F_t(x)$ be the \emph{frontier}.  It contains the candidates with the
highest cached score among those for which $x\in\Omega_t(s)$, including all
ties.
In the all-failure example above, every candidate is on the frontier and each
receives zero per-instance credit.  Define the legally observed comparison
set and the number of candidates on the frontier by
\[
\mathcal C_t(x)=\{u\mid x\in\Omega_t(u)\},\qquad
N_x=|\mathcal C_t(x)|,\qquad
y_{sx}=\ind\{s\in\F_t(x)\},\qquad
K_x=\sum_{u\in\mathcal C_t(x)}y_{ux}.
\]
When $N_x\ge2$, candidate $s$ receives the per-instance credit
\[
r_{sx}=y_{sx}-\frac{K_x-y_{sx}}{N_x-1}
=\frac{N_xy_{sx}-K_x}{N_x-1}.
\]
Only comparable exposures enter the aggregate.
\begin{equation}
\begin{aligned}
\Omega_t^{\mathrm{cmp}}(s)&=\{x\in\Omega_t(s)\mid N_x\ge2\},
&E_s^{\mathrm{cmp}}&=|\Omega_t^{\mathrm{cmp}}(s)|,\\
R_s&=\sum_{x\in\Omega_t^{\mathrm{cmp}}(s)}r_{sx},
&G_s&=\frac{R_s}{E_s^{\mathrm{cmp}}}.
\end{aligned}
\label{eq:global_key}
\end{equation}
The relative credit $G_s$ is the average per-instance credit.  It records how
often $s$ is on the current frontier relative to the
other legally observed candidates on the same instances.  An observation with
$N_x=1$ remains in the committed table but contributes no relative credit until
another candidate is legally observed on $x$.  Missing entries remain missing,
and $G_s$ is undefined when $E_s^{\mathrm{cmp}}=0$ rather than being set to
zero.  Whenever a commit changes an observation or a frontier, COMPASS
recomputes $N_x$, $K_x$, all affected per-instance credits, and the resulting $G_s$ values.

\subsection{A credit-ranked active set}
\label{sec:lineage_compression}

At each wave, COMPASS considers all accepted candidates whose relative credit is
defined.  It keeps the highest-credit candidates up to the active-rank cutoff
$\rho$, including every candidate tied at the cutoff.  Parents and descendants
are considered together.  If fewer than $\rho$ candidates have defined credit,
all of them remain active.  The tie rule can therefore leave more than $\rho$
active programs, and a parent and its descendant can remain active at the same
time.  If no candidate has defined credit, COMPASS uses the seed.

Proposal parents are selected from the active set in decreasing $G_s$ order,
with candidate index breaking exact ties.  The order is traversed cyclically
across waves, and each proposal slot defines one revision attempt.

\subsection{Turning credit into proposal work and legal admission}
\label{sec:proposal_allocation}
\label{sec:admission_allocation}
\label{sec:commit}

Relative credit $G_s$ selects which candidates to revise, and matching determines
which proposal batches they receive.  COMPASS draws $w$ proposal batches and
matches them one-to-one to the $w$ proposal slots.  Proposal observations
remain available for diagnosis, so matching uses the full cache, including
observations excluded from the parent's credit.  The matching weight $H(s,P)$
averages $s$'s gap to the highest cached score among active candidates,
computed separately for each instance in $P$.
\begin{equation}
H(s,P)=\frac{1}{|P|}\sum_{x\in P}
\begin{cases}
\left[\max_{\substack{r\in\A_t\\q_t(r,x)\ \mathrm{finite}}}q_t(r,x)-q_t(s,x)\right]_+,
& q_t(s,x)\ \mathrm{finite},\\
0,&\text{otherwise}.
\end{cases}
\label{eq:repairable_gap}
\end{equation}
Here $[v]_+=\max\{v,0\}$.  Missing parent observations contribute zero to this
matching weight and remain missing in the cache.  Let
$\mathcal I_w=\{0,\ldots,w-1\}$ index the drawn slots $z_i$ and proposal
minibatches $P_j$.  \method{} solves the maximum-weight matching, written as the binary
program
\begin{equation}
\begin{aligned}
Y^\star\in\argmax_{Y\in\{0,1\}^{w\times w}}\quad
&\sum_{i,j\in\mathcal I_w}H(z_i,P_j)y_{ij}\\
\text{s.t.}\quad
&\sum_{j\in\mathcal I_w}y_{ij}=1\ \ \forall i,
\qquad
\sum_{i\in\mathcal I_w}y_{ij}=1\ \ \forall j.
\end{aligned}
\label{eq:proposal_assignment}
\end{equation}
Thus distinct slots receive distinct proposal batches even when two slots
carry the same parent label.  The assignment maximizes the sum of their
matching weights.

Proposal evaluation updates the existing parent's cached observations.
Tasks with usable execution traces and reflection data proceed to child
generation.  When perfect-score skipping is enabled, a task also requires at
least one proposal score below the configured maximum.  Write $m\le w$ for
the number of these remaining tasks, with parent $s_i$ and proposal batch $P_i$.
For each task independently, the benchmark evaluator samples an admission batch
$A_i$ satisfying
\[
A_i\cap X(s_i,P_i)=\varnothing.
\]
The evaluator freezes $A_i$ and, for
each $x\in A_i$, selects a reference from the current frontier with
defined $G$, preferring the largest $G$ and using seeded randomness for an
exact tie.  If no frontier candidate has defined credit, the sampled parent $s_i$ is the
reference.  Denote this choice by $r_i(x)$.  Missing reference observations are
evaluated before the child is generated and stored using the highest-score
cache rule.  The admission instances, references, and outcomes remain hidden
from the generator.

After the generator fixes child $p_i$ from $(s_i,P_i)$, the evaluator scores it
on $A_i$ and compares its scores with the frozen reference scores.  In this
gate, $q(p_i,x)$ is the child's new score and $q(r_i(x),x)$ is the fixed reference
score on the same instance.  The paired aggregate gate is
\begin{equation}
\Delta_i=\sum_{x\in A_i}\left[q(p_i,x)-q(r_i(x),x)\right],
\qquad
\operatorname{admit}(p_i)\iff\Delta_i>0.
\label{eq:admission_gate}
\end{equation}
Proposal and reference evaluations update existing programs' cached scores.
Only credit-eligible entries affect their credit.  An accepted child's
admission observations enter the cache with that child.  Its proposal batch
joins the inherited proposal history, and later observations outside that set can
also contribute to its credit.
COMPASS then recomputes every affected $G_s$, the tie-inclusive active set,
and the queue.  At the budget boundary, COMPASS selects the program with the
largest defined $G_s$, with earliest index breaking an exact tie.  The selected
program then receives independent benchmark evaluation.  Expanded execution
details appear in \Cref{app:algorithm}.

\section{How \method{} Turns Evidence into Search}
\label{sec:mechanism}

To see how the components fit together, follow one revision through the search
loop.  A proposal batch reveals a parent's failures and shapes a child.  Since
those examples helped create the child, different evidence must establish that
the revision is useful.  The admission outcomes are compared with other
candidates on the same instances to update their relative credit, which selects
parents for the next proposal attempts.  The formal statements below support
three intuitions needed to read this loop.  Admission tests a fixed revision,
credit records relative performance on shared instances, and a fixed
rollout budget supports both candidate generation and evaluation.

\subsection{A revision is judged by paired evaluation}

Admission asks whether a fixed revision improves the total score on its
evaluation batch.  The child and its references are compared on the same
 instances, and the child is admitted when the total paired difference is
 positive.

\begin{proposition}[Reusing proposal examples can overstate performance]
\label{prop:self_cert_leakage}
For every admission size $n\ge1$ and $\epsilon>0$, there is a task
distribution and a proposal rule that observes $n$ examples, constructs a
program, and scores one when those examples are reused for evaluation,
although the program's population utility is at most $\epsilon$.
\end{proposition}

A descendant can inherit behavior adapted to an ancestor's proposal examples.
COMPASS therefore carries the exclusion forward along the lineage.  The
construction in \Cref{app:proofs} illustrates this with memorization.

\paragraph{Paired evaluation.}
The proposal batch guides revision of the parent.  The evaluator samples the
admission batch outside $X(s,P)$, fixes its per-instance references, and hides
this plan from the generator.  Once the child is fixed, the gate in
\Cref{eq:admission_gate} compares it with those references on the same
instances.  \Cref{thm:paired_admission} establishes lineage exclusion and
bounds the deviation of the mean paired difference from the average
conditional mean.

\subsection{Admission outcomes become relative credit}

Admission decides whether a particular child enters the pool.  Routing asks
which accepted candidates have shown distinctions worth pursuing further.
COMPASS answers it by comparing candidates on the same instance before
aggregating across the partial cache.

\begin{proposition}[Instance-wise credit conservation]
\label{prop:credit_conservation}
For every instance $x$ with $N_x\ge2$, the per-instance credits defined immediately
before \Cref{eq:global_key} satisfy
\[
\sum_{s\in\mathcal C_t(x)}r_{sx}=0.
\]
\end{proposition}

This conservation law assigns credit to differences between candidates.  If all
$N_x$ compared candidates are on the frontier, the instance assigns
zero per-instance credit to every candidate.  If one
candidate is the only one on the frontier, it receives credit of one and each of the
other $N_x-1$ candidates receives $-1/(N_x-1)$.  Thus being on the frontier
receives more credit when it distinguishes a candidate from its observed peers.

New results can change relative credit even when the cached scores of existing
candidates do not change.  When another candidate becomes comparable on an
instance, COMPASS recomputes that instance's frontier and credits.  An instance
with one legal observation begins contributing relative credit when a second
candidate becomes comparable there.

\subsection{Credit directs the next wave}

Matching accounts for parents competing for the same proposal batch.

For illustration, suppose two parent slots have the following matching weights.
\[
\begin{array}{c|cc}
 & P_1 & P_2\\ \hline
a & 0.40 & 0.30\\
b & 0.35 & 0.05
\end{array}
\]
Both prefer $P_1$.  Assigning $P_1$ to $a$ leaves $P_2$ for $b$, for a total
weight of $0.45$.  Assigning $P_2$ to $a$ and $P_1$ to $b$ gives $0.65$.
Maximum-weight matching accounts for these competing preferences across the
whole wave.  After paired admission, accepted outcomes update the credit used
in the next wave.
At the budget boundary, the program with maximal defined $G_s$ is returned.

\subsection{Search and final selection}
\label{sec:search_selection}

A search can miss a strong program or discover one and return a weaker
candidate.  The search gap and selection gap separate these two sources of
regret.  Let $\Pcal_B$ be the candidates available when search stops and let
$\hat p_B$ be the program selected by the largest defined $G_s$.  Its regret
has the exact decomposition
\begin{equation}
R^\star-R(\hat p_B)
=
\underbrace{R^\star-\max_{p\in\Pcal_B}R(p)}_{\text{search gap}}
+
\underbrace{\max_{p\in\Pcal_B}R(p)-R(\hat p_B)}_{\text{selection gap}}.
\label{eq:terminal_regret_decomposition}
\end{equation}
Credit-ranked slots and proposal matching allocate candidate-generation work.
paired admission determines which revisions enter the pool.
To relate selection by relative credit to selection by sample mean,
\Cref{prop:bandwidth_precision} first bounds the error of mean-based selection.
\Cref{cor:g_selection} then connects it to the candidate chosen by $G$.
Consider $M$ candidates,
each fixed before receiving $n$ conditionally independent and identically
distributed rewards under a common benchmark distribution, with mean $R(p_i)$.
Let $\widehat R_n(p_i)$ be the arithmetic mean of these rewards and let
$\widehat p_G$ maximize defined credit in this pool.  Define
\begin{equation}
\xi_G=\max_{i\le M}\widehat R_n(p_i)-\widehat R_n(\widehat p_G).
\label{eq:g_selection_gap}
\end{equation}
Under these evaluation conditions, \Cref{cor:g_selection} gives, with
probability at least $1-\delta$,
\begin{equation}
R^\star-R(\widehat p_G)
\le \underbrace{R^\star-\max_{i\le M}R(p_i)}_{\text{search gap}}
+\underbrace{2\sqrt{\frac{\log(2M/\delta)}{2n}}+\xi_G}_{\text{selection gap bound}}.
\label{eq:g_selection_regret}
\end{equation}
The term $\xi_G$ is the difference in sample mean between the highest-mean
candidate and the candidate selected by $G$.  With complete shared binary observations and one score
per candidate--instance pair, the two rankings coincide and $\xi_G=0$
(\Cref{eq:binary_credit_mean}).

\section{Experiments}
\label{sec:experiments}

\subsection{Evaluation questions}

The experiments evaluate whether provenance-constrained routing improves held-out benchmark
utility under a fixed budget and whether the result persists across models and
tasks.

\subsection{Tasks, models, and evaluation}

The study spans AIME-2025 and LiveBench-Math mathematical reasoning, IFBench
instruction following with unseen constraint types, HotPotQA question
answering, HoVer multi-hop fact verification, PUPA privacy-conscious
delegation
\citep{ifbench_bench,aime_2025,livebench,hotpotqa_bench,hover_bench,
papillon_bench}.  The models are Qwen3-8B~\citep{qwen3_technical_report} and
GPT-5.4-nano.  PUPA retains its official GPT-4.1-mini judge for response quality
and leakage of personally identifiable information (PII).  This judge remains
fixed throughout optimization and testing.

Official evaluation is separated from optimization.  Optimization uses the
per-task rollout budgets reported for GEPA~\citep{gepa}.  The budgets are 6,871
for HotPotQA, 3,593 for IFBench, 7,051 for HoVer, 2,426 for PUPA, 1,839 for
AIME-2025, and 1,839 for LiveBench-Math.  The Qwen3-8B AIME variant is
evaluated on all 150 benchmark problems.  HoVer uses 150/300/300 train,
validation, and test instances.  LiveBench-Math uses its frozen 126-problem
test.  Candidate selection uses optimization data, followed by one evaluation
on the held-out test split.

\subsection{Comparison across six tasks with Qwen3-8B}

Table~\ref{tab:qwen_results} compares Qwen3-8B evaluations on all six shared
tasks.  Seed, GRPO, MIPROv2, GEPA, and GEPA+Merge reuse published GEPA-study
task scores.  We compare two COMPASS ablations.  \mbox{\emph{w/o Credit \& Matching}}
selects top-ranked candidates using frontier frequency over legal exposures
($F/E$) and omits skill-proposal matching.
\mbox{\emph{w/o Matching}} retains ranking by $G_s$ in \Cref{eq:global_key} and omits
only matching in \Cref{eq:proposal_assignment}.
The COMPASS row reports the program selected for each task.  COMPASS
has the highest six-task aggregate and the highest score on all six tasks,
including PUPA, where it exceeds GEPA by $1.19$ percentage points.

\begin{table}[H]
\caption{\textbf{Qwen3-8B evaluation across all shared tasks (\%).}
Each cell reports the official task aggregate.  Aggregate is the unweighted
mean over the six tasks, and $\Delta$ Seed is the aggregate improvement in
percentage points.  Both are computed from the displayed task scores before
rounding to two decimal places.  The first five rows reuse GEPA-study task
scores.  Ablation names indicate the removed components.  Credit denotes
$G_s$ ranking, and Matching denotes skill-proposal matching.
TextGrad and Trace were not reported for Qwen3-8B.}
\label{tab:qwen_results}
\begin{center}
\setlength{\tabcolsep}{3.3pt}
\begin{tabular*}{\linewidth}{@{\extracolsep{\fill}}lrrrrrrrr@{}}
\toprule
Method & HotPotQA & IFBench & HoVer & PUPA & \shortstack[r]{AIME-\\2025} & \shortstack[r]{LiveBench-\\Math} & Aggregate & $\Delta$ Seed \\
\midrule
Seed & 42.33 & 36.90 & 35.33 & 80.82 & 27.33 & 48.70 & 45.24 & -- \\
GRPO & 43.33 & 35.88 & 38.67 & 86.66 & 38.00 & 51.26 & 48.97 & +3.73 \\
MIPROv2 & 55.33 & 36.22 & 47.33 & 81.55 & 20.00 & 46.60 & 47.84 & +2.60 \\
GEPA & 62.33 & 38.61 & 52.33 & 91.85 & 32.00 & 51.95 & 54.85 & +9.61 \\
GEPA+Merge & 64.33 & 28.23 & 51.67 & 86.26 & 32.00 & 51.95 & 52.41 & +7.17 \\
\begin{tabular}[c]{@{}l@{}}w/o~Credit \\ \& Matching\end{tabular} & 64.67 & 40.31 & 51.67 & 86.91 & 57.33 & 66.67 & 61.26 & +16.03 \\
w/o~Matching & 65.33 & 40.82 & 50.67 & 88.85 & 60.67 & 67.68 & 62.34 & +17.10 \\
\method{} & \textbf{67.00} & \textbf{43.37} & \textbf{54.67} & \textbf{93.04} & \textbf{63.33} & \textbf{73.96} & \textbf{65.90} & \textbf{+20.66} \\
\bottomrule
\end{tabular*}
\end{center}
\end{table}

The ablations separate the effects of credit ranking and matching.
\mbox{\emph{w/o Credit \& Matching}}, \mbox{\emph{w/o Matching}}, and COMPASS attain aggregate
scores of $61.26\%$, $62.34\%$, and $65.90\%$, respectively.  Comparing the two
ablations evaluates the ranking signal, while comparing \mbox{\emph{w/o Matching}} with
COMPASS evaluates the addition of skill-proposal matching.

The six tests contain 300, 294, 300, 221, 150, and 126 instances for HotPotQA,
IFBench, HoVer, PUPA, AIME-2025, and LiveBench-Math, respectively.  Candidate
selection precedes the official test.  The selected-candidate records are
listed in \Cref{app:protocol}.

\subsection{Same-model comparison with GPT-5.4-nano}

Table~\ref{tab:nano_results} reports the GPT-5.4-nano comparison.  Both task
execution and reflection use GPT-5.4-nano.  PUPA keeps the fixed
benchmark scorer described above.  COMPASS has the highest six-task aggregate
and the highest score on four tasks.  MIPROv2 is higher on HotPotQA, and GEPA is
higher on HoVer.

\begin{table}[H]
\caption{\textbf{GPT-5.4-nano evaluation across all shared tasks (\%).}
Each cell reports the official task aggregate.  Aggregate is the unweighted
mean over the six tasks, and $\Delta$ Seed is the aggregate improvement in
percentage points.  The table includes methods with completed same-model
evaluations.}
\label{tab:nano_results}
\begin{center}
\setlength{\tabcolsep}{3.3pt}
\begin{tabular*}{\linewidth}{@{\extracolsep{\fill}}lrrrrrrrr@{}}
\toprule
Method & HotPotQA & IFBench & HoVer & PUPA & \shortstack[r]{AIME-\\2025} & \shortstack[r]{LiveBench-\\Math} & Aggregate & $\Delta$ Seed \\
\midrule
Seed & 31.33 & 48.98 & 45.33 & 74.86 & 15.33 & 35.09 & 41.82 & -- \\
MIPROv2 & \textbf{61.67} & 47.79 & 53.00 & 75.45 & 30.00 & 43.66 & 51.93 & +10.11 \\
GEPA & 57.33 & 48.30 & \textbf{61.33} & 81.09 & 34.00 & 35.03 & 52.85 & +11.03 \\
\method{} & 60.00 & \textbf{51.36} & 60.33 & \textbf{84.14} & \textbf{42.00} & \textbf{45.03} & \textbf{57.14} & \textbf{+15.32} \\
\bottomrule
\end{tabular*}
\end{center}
\end{table}

\section{Related Work}
\label{sec:related}

\paragraph{Reflective textual-program optimization.}
Reflexion and Self-Refine use language feedback to revise behavior, while
TextGrad, Trace/OptoPrime, MIPROv2, and GEPA optimize textual components from
execution evidence
\citep{reflexion,self_refine,textgrad,
cheng2024traceautodiffgenerativeoptimization,miprov2,gepa}.  COMPASS uses the
resulting scores to choose which candidates to revise next and which batches
to assign them.

\paragraph{Evolution, illumination, and skill search.}
Evolutionary prompt and skill methods maintain multiple candidates to avoid a
single greedy trajectory~\citep{gepa,skillpro,grasp,evoskills,metaskillevolve}.
COMPASS keeps parents and descendants in one credit-ranked candidate pool.
Its active set contains the highest-credit candidates, including all ties at
the cutoff.

\paragraph{Adaptive evaluation and budget allocation.}
TRIPLE, POES, and Bayesian-Agent allocate evaluation or retrieval effort under
limited budgets~\citep{triple,poes,bayesianagent}. TextBO and APEX optimize
evaluation efficiency and development-example choice
\citep{kang2026textbo,wang2026apex}.  COMPASS allocates proposal attempts by
relative credit and assigns their batches jointly by maximum-weight matching.
Both decisions use the existing execution records.

\paragraph{Generalization after adaptive search.}
Repeated adaptive use of optimization data can widen the gap between selected
validation performance and held-out test performance
\citep{dwork2015generalization}.  GEPA addresses this empirically through
reusable rules and branch diversity~\citep{gepa}.  COMPASS carries proposal
history through the lineage.  Examples used to create an ancestor remain
excluded from a descendant's admission.  Held-out benchmark tests evaluate
the selected program.

\section{Conclusion}

COMPASS uses execution results to decide which candidates to develop and which
batches they should work on.  Proposal history determines the eligible
comparisons, and new results update the relative credit of candidates already
in the pool.  Relative credit selects proposal parents, and maximum-weight
matching assigns their batches.  Paired admission evaluates the resulting
revisions on separate examples.  Under the reported budgets, COMPASS achieves the highest six-task
aggregate for both Qwen3-8B and GPT-5.4-nano, and leads on all six Qwen3-8B tasks.
